\subsection{Coordinate limits and integral proper maps}
\label{subsec:spectrum-filters}

\begin{lemma}[Ultrafilter limits in closed evaluation spaces]
\label{spectrum:closed-coordinate-ultrafilter-limit}
\lean{EnvelopeProperty.exists_ultrafilter_limit_of_bounded_coordinates,
EnvelopeProperty.Berkovich.Spectrum.exists_ultrafilter_limit_of_bounded_evaluations,
EnvelopeProperty.Berkovich.RelativeSpectrum.exists_ultrafilter_limit_of_bounded_evaluations}
\leanok
\uses{spectrum:topology,spectrum:closed-evaluation-embedding,spectrum:relative}
Let $X$ be a topological space with a closed embedding into a product
of copies of $\mathbb R_{\ge0}$. If an ultrafilter on $X$ is eventually
bounded above in each coordinate, then it converges to a point of $X$.
In particular this holds for absolute seminorm spectra and for spectra
with any prescribed base seminorm. The coordinate bounds and their
eventual sets need not be uniform across the coordinates.
\end{lemma}
\begin{proof}
For each coordinate, its mapped ultrafilter contains a compact interval
$[0,M]$, hence converges to some value in that interval. Choose these
values coordinate by coordinate. Convergence in the product topology
gives a limit of the whole evaluation-function ultrafilter. Closedness
of the embedding's image places this limit in that image; choose its
preimage. An embedding reflects convergence, so the original ultrafilter
converges to that preimage. The absolute seminorm evaluation embedding
is closed, and its composition with the closed prescribed-base inclusion
gives the relative case.
\end{proof}

\begin{theorem}[Integral pullbacks are proper and closed]
\label{spectrum:integral-proper-map}
\lean{EnvelopeProperty.Berkovich.Spectrum.isProperMap_comap_of_isIntegral,
EnvelopeProperty.Berkovich.Spectrum.isClosedMap_comap_of_isIntegral,
EnvelopeProperty.Berkovich.RelativeSpectrum.isProperMap_comap_of_isIntegral,
EnvelopeProperty.Berkovich.RelativeSpectrum.isClosedMap_comap_of_isIntegral}
\leanok
\uses{spectrum:closed-coordinate-ultrafilter-limit,spectrum:integral-root-bound,
spectrum:pullback,spectrum:relative}
For an integral homomorphism $\varphi:A\to B$ of arbitrary commutative
rings, pullback $\mathcal S(B)\to\mathcal S(A)$ is a proper continuous
map, and in particular is closed. For an integral algebra homomorphism
over an arbitrary ring with a prescribed real seminorm, the relative
pullback map is proper and closed as well. These assertions require
no finite-generation or local-compactness hypothesis.
\end{theorem}
\begin{proof}
Suppose an ultrafilter in $\mathcal S(B)$ has convergent image in
$\mathcal S(A)$. Every $A$-evaluation of that image is eventually bounded.
For each $b\in B$, choose a monic integral equation. Finitely many
coefficient evaluations admit simultaneous eventual bounds; their finite
maximum, enlarged to at least one, bounds $|b|$ by the monic-root estimate.
This gives an eventual bound for each $B$-evaluation, with no uniform
bound over all $b$ required. The preceding lemma supplies a source limit.
Continuity and the Hausdorff target identify its image with the prescribed
target limit. This is the ultrafilter criterion for properness, which
also implies closedness.
For the relative map, compose the closed prescribed-base inclusion into
$\mathcal S(B)$ with the absolute proper map. Closed embeddings are
proper, so this composition is proper. It equals relative pullback
followed by the injective prescribed-base inclusion into $\mathcal S(A)$.
Properness of a composition with an injective continuous outer map
implies properness of the inner map. Applying this cancellation gives
relative properness and hence closedness.
This argument proves genuine filter lifting rather than inferring
properness from compact inverse images alone.
\end{proof}
